%LTeX: language=fr-FR
\documentclass[12pt]{article}
\usepackage[margin=0.8in]{geometry}
\usepackage[utf8]{inputenc}
\usepackage{graphicx}
\usepackage[export]{adjustbox}
\usepackage{amsmath}
\usepackage{amsfonts}
\usepackage{amssymb}
\usepackage[version=4]{mhchem}
\usepackage{stmaryrd}
\usepackage{bbold}
%\usepackage[fallback]{xeCJK}
%\usepackage{polyglossia}
\usepackage{fontspec}
\usepackage{censor}
\usepackage{currfile}
%\setCJKmainfont{Noto Serif CJK TC}
\usepackage{enumitem}
\usepackage{tikz}
\usetikzlibrary {arrows.meta, quotes} 
\usepackage{tkz-tab}
\usepackage{pgfplots}
\pgfplotsset{compat=1.18}
\usepackage{amsthm}
\usepackage{tcolorbox}
\usepackage{array}
\usepackage{colortbl}
\usepackage{tabularray}
\usepackage{tabularx}
\usepackage{makecell}
\usepackage{esvect}
%\usepackage{changepage}   % for the adjustwidth environment
\usepackage{enumitem}
\usepackage{tasks}
\usepackage{fancyhdr}
\usepackage[french]{babel}
\usepackage{lastpage}

\tcbuselibrary{theorems}

\definecolor{GrayCen}{HTML}{d2d2d2}
\definecolor{GrayBG}{HTML}{d9d9d9}
\definecolor{BlueDef}{HTML}{104e8b}
\definecolor{GreenProp}{HTML}{025540}
\definecolor{RedMet}{HTML}{cd5556}
\definecolor{BlackSav}{HTML}{000000}
\def\censorcolor{GrayCen}
\let\svcensorrule\censorrule
\renewcommand\censorrule[1]{%
\textcolor{\censorcolor}{\svcensorrule{#1}}}

\newtcbtheorem
[]% init options
{definition}% name
{Définition}% title
{%
  colback=GrayBG,
  colframe=BlueDef,
  fonttitle=\bfseries,
}% options
{def}% prefix

\newtcbtheorem
[]% init options
{proposition}% name
{Proposition}% title
{%
  colback=GrayBG,
  colframe=GreenProp,
  fonttitle=\bfseries,
}% options
{prop}% prefix

\newtcbtheorem
[]% init options
{method}% name
{Méthode}% title
{%
  theorem name,
  colback=GrayBG,
  colframe=RedMet,
  fonttitle=\bfseries,
}% options
{met}% prefix

\newtcbtheorem
[]% init options
{savoir}% name
{Savoir-faire du chapitre}% title
{%
  theorem name,
  colback=GrayBG,
  colframe=BlackSav,
  fonttitle=\bfseries,
}% options
{sav}% prefix

\newtheoremstyle{break}
{\topsep}{\topsep}%
{}{}%
{\bfseries}{}%
{\newline}{}%
\theoremstyle{break}
\newtheorem*{example}{Exemple~:}
\newtheorem*{examples}{Exemples~:}
\newtheorem*{remark}{Remarque~:}
\newtheorem*{remarks}{Remarques~:}
% Variante sans retour à la ligne après « Exemple. »
\newtheoremstyle{inline}
{\topsep}{\topsep}%
{}{}%
{\bfseries}{}%
{ }{}%
\theoremstyle{inline}
\newtheorem*{example*}{Exemple~:}

\newenvironment{demo}[1][\proofname]{%
  \begin{proof}[#1]$ $\newline
  }{%
  \end{proof}
}

%\setmainfont{CMU Serif}name

\title{
  \vspace{4em}
  \centering
  \tcbset{colframe=BlackSav,colback=GrayBG}
  \tcbox[tikznode]{
    Chapitre 4 \\ Généralités sur les suites
  }
}

\author{}
\date{}

\pagestyle{empty}
\graphicspath{ {../Images/} }
\input{\currfiledir ../def.tex}
\def \Document {Cours}


\setlength\parindent{0pt}
%\setlist[itemize]{topsep=0pt}
\setlist{nolistsep}
%\setlist[itemize]{noitemsep}
%\setlist[enumerate]{noitemsep}

\tikzset{cross/.style={path picture={
  \draw[black]
    (path picture bounding box.south east)--(path picture bounding box.north west)
    (path picture bounding box.south west)--(path picture bounding box.north east);
}}}

\newcolumntype{C}[1]{>{\centering\arraybackslash}m{#1}}
\newcolumntype{L}[1]{>{\raggedright\arraybackslash}m{#1}}
\newcolumntype{N}{@{}m{0pt}@{}}

\newcolumntype{Y}{>{\centering\arraybackslash}X}

\fancypagestyle{firststyle}
{
  \fancyhf{}
  \fancyhead[L]{\Structure}
  \fancyhead[R]{\Niveau}
  \fancyfoot[L]{Alexandre Fernandez}
  \fancyfoot[R]{{\thepage\  / \pageref{LastPage}}}
   %\renewcommand{\headrulewidth}{0pt} % removes horizontal header line
}


\begin{document}

\maketitle
\thispagestyle{firststyle}

\tableofcontents

% \newcommand{\multihide}[1]{\mbox{\vphantom{#1}\colorbox{gray!50}{\parbox{\textwidth}{{#1}}}}}
%\newcommand{\multihide}[1]{\mbox{\vphantom{#1}\colorbox{gray!50}{\parbox{\textwidth}{{\phantom{#1}}}}}}
% On crée une boîte temporaire pour stocker le texte
\newsavebox{\hidebox}

% -----------
% Macros des trous : version élève si \ELEVE est défini à la compilation
% (recette « prof + élève »), sinon version prof.
\ifdefined\ELEVE
% Version élève : contenu invisible
\newcommand{\hideM}[1]{\mbox{\vphantom{$#1$}\smash{\colorbox{gray!50}{\phantom{$\displaystyle #1$}}}} }
\newcommand{\hideMW}[1]{\mbox{\vphantom{$#1$}\smash{\colorbox{white}{\phantom{$\displaystyle #1$}}}} }
\newcommand{\hide}[1]{\mbox{\vphantom{#1}\smash{\colorbox{gray!50}{\phantom{#1}}}}}
\NewEnviron{multihide}{%
  \colorbox{gray!50}{
    \begin{lrbox}{\hidebox}%
        \begin{minipage}{\linewidth}%
          \BODY
        \end{minipage}%
    \end{lrbox}%
    % On affiche un espace invisible ayant les dimensions exactes de la boîte
    \vrule width 0pt height \ht\hidebox depth \dp\hidebox\hskip\wd\hidebox%
  }
}
% Corrigé long dans un cadre blanc : zone vide de même hauteur
% Argument optionnel : espace ajouté sous le corrigé, ex. \begin{multihideW}[2em]
\NewEnviron{multihideW}[1][0pt]{%
  \par\noindent\fbox{%
    \begin{lrbox}{\hidebox}%
      \begin{minipage}{\dimexpr\linewidth-2\fboxsep-2\fboxrule}%
        \BODY\par\vspace{#1}%
      \end{minipage}%
    \end{lrbox}%
    \vrule width 0pt height \ht\hidebox depth \dp\hidebox\hskip\wd\hidebox%
  }
}
\newcommand{\hidetk}[1]{}
\else
% Version prof : contenu visible sur fond gris
\newcommand{\hideM}[1]{\mbox{\vphantom{$#1$}\smash{\colorbox{gray!50}{$\displaystyle #1$}}} }
\newcommand{\hideMW}[1]{\mbox{\vphantom{$#1$}\smash{\colorbox{white}{$\displaystyle #1$}}} }
\newcommand{\hide}[1]{\mbox{\vphantom{#1}\smash{\colorbox{gray!50}{#1}}}}
\NewEnviron{multihide}{%
  \colorbox{gray!50}{
    \begin{minipage}{\linewidth}%
      \BODY
    \end{minipage}%
  }
}
% Corrigé long dans un cadre blanc (argument optionnel : espace ajouté)
\NewEnviron{multihideW}[1][0pt]{%
  \par\noindent\fbox{%
    \begin{minipage}{\dimexpr\linewidth-2\fboxsep-2\fboxrule}%
      \BODY\par\vspace{#1}%
    \end{minipage}%
  }
}
\newcommand{\hidetk}[1]{#1}
\fi
% -----------

\newcommand{\wvec}[1]{\overrightarrow{#1}}
\newcommand{\ssi}[0]{\emph{ssi}\ }
\newcommand{\setR}[0]{$\mathbb{R}$}

\newpage

\pagestyle{fancy}

\fancyhf{} % clear existing header/footer entries
\fancyhead[R]{ {\Niveau}}
\fancyhead[L]{ {\Chapitre}}
%\fancyhead[R]{\large \textbf{\textbf{\Document\ (B) -- 1h}}}
\fancyfoot[L]{ A. Fernandez }
\fancyfoot[R]{{\thepage\  / \pageref{LastPage}}}

% Repère quadrillé : \repere{xmin}{xmax}{ymin}{ymax}
\newcommand{\repere}[4]{%
  \draw[very thin, gray!40, step=1] (#1,#3) grid (#2,#4);
  \draw[->] (#1,0) -- ({#2+0.4},0) node[below right] {\small $x$};
  \draw[->] (0,#3) -- (0,{#4+0.4}) node[above left] {\small $y$};
  \foreach \x in {#1,...,#2} {\ifnum\x=0\else\draw (\x,0.08) -- (\x,-0.08) node[below] {\scriptsize $\x$};\fi}
  \foreach \y in {#3,...,#4} {\ifnum\y=0\else\draw (0.08,\y) -- (-0.08,\y) node[left] {\scriptsize $\y$};\fi}
  \node[below left] at (0,0) {\scriptsize $0$};
}

% Cadre pour les programmes Python
\newtcolorbox{python}{colback=white, colframe=black!50, boxrule=0.5pt, arc=1pt, left=6pt, fontupper=\ttfamily}
\newcommand{\tab}{\hspace*{2em}}

\section{Notion de suite et modes de génération}

\subsection{Définition et notations}

\begin{example}
  Dans l'activité 1, on a noté $u(n)$ le nombre d'enfants inscrits à la crèche à la rentrée $2026 + n$.

  \smallskip
  \begin{minipage}[c]{0.52\linewidth}
    \renewcommand{\arraystretch}{1.4}
    \begin{tabular}{|c|C{1cm}|C{1cm}|C{1cm}|C{1cm}|C{1cm}|}
      \hline
      $n$ & $0$ & $1$ & $2$ & $3$ & $4$ \\
      \hline
      $u(n)$ & $20$ & $26$ & $32$ & $38$ & $44$ \\
      \hline
    \end{tabular}
  \end{minipage}
  \hfill
  \begin{minipage}[c]{0.45\linewidth}
    Ces nombres forment une liste numérotée par $0$, $1$, $2$, \dots{} : c'est une \textbf{suite}.
  \end{minipage}
\end{example}

\begin{definition}{Suite numérique}{}
  \begin{multihide}
    Une \textbf{suite} $u$ associe à chaque entier naturel $n$, appelé \textbf{rang}, un nombre $u(n)$, appelé \textbf{terme de rang $n$}. Le \textbf{premier terme} est en général $u(0)$ (parfois $u(1)$).
  \end{multihide}
\end{definition}

\begin{remark}
  Le rang est un entier naturel : $n \in \mathbb{N} = \{0~;~1~;~2~;~3~;~\dots\}$. Ainsi, $u(2{,}5)$ n'existe pas.
\end{remark}

\begin{definition}{Notation $u_n$}{}
  \begin{multihide}
    Le terme $u(n)$ se note aussi $u_n$ (« $u$ indice $n$ ») : $u_0 = u(0)$, $u_{n+1} = u(n+1)$, etc.

    Dans la suite du cours, on utilise la notation $u_n$.
  \end{multihide}
\end{definition}

\begin{example}
  Une suite $u$ est donnée par le tableau ci-dessous.

  \smallskip
  \begin{minipage}[c]{0.52\linewidth}
    \renewcommand{\arraystretch}{1.4}
    \begin{tabular}{|c|C{1cm}|C{1cm}|C{1cm}|C{1cm}|C{1cm}|}
      \hline
      $n$ & $0$ & $1$ & $2$ & $3$ & $4$ \\
      \hline
      $u_n$ & $5$ & $8$ & $4$ & $10$ & $7$ \\
      \hline
    \end{tabular}
  \end{minipage}
  \hfill
  \begin{minipage}[c]{0.45\linewidth}
    \begin{itemize}
      \item Le premier terme est $u_0 = \hideM{5}$.
      \item Le terme de rang $2$ est $u_2 = \hideM{4}$.
      \item Le terme $10$ a pour rang \hideM{3}.
    \end{itemize}
  \end{minipage}
\end{example}

\begin{remark}
  Attention : $u_{n+1}$ est le terme de rang $n+1$ (le terme \emph{suivant}). Ce n'est pas $u_n + 1$.

  Par exemple, dans le tableau ci-dessus : $u_{2+1} = u_3 = \hideM{10}$, alors que $u_2 + 1 = \hideM{4 + 1 = 5}$.
\end{remark}


\subsection{Formule explicite}

\begin{example}
  Dans l'activité 1, $u_n = 20 + 6n$. À la rentrée 2036 ($n = 10$) : $u_{10} = \hideM{20 + 6 \times 10 = 80}$ enfants.
\end{example}

\begin{definition}{Formule explicite}{}
  \begin{multihide}
    Une suite est définie par une \textbf{formule explicite} quand $u_n$ est donné \textbf{en fonction de $n$}.

    On calcule directement n'importe quel terme, comme une image par une fonction.
  \end{multihide}
\end{definition}

\begin{method}{Calculer un terme avec une formule explicite}{}
  Pour calculer un terme, on \textbf{remplace $n$} par son rang dans la formule.
\end{method}

\begin{examples}
  Calculer le premier terme et le terme de rang $5$ de chaque suite.

  \medskip
  \begin{minipage}[t]{0.31\linewidth}
    $u_n = 3n + 2$
    \begin{multihideW}[1em]
      $u_0 = 3 \times 0 + 2 = 2$

      $u_5 = 3 \times 5 + 2 = 17$
    \end{multihideW}
  \end{minipage}
  \hfill
  \begin{minipage}[t]{0.31\linewidth}
    $v_n = n^2 + 1$
    \begin{multihideW}[1em]
      $v_0 = 0^2 + 1 = 1$

      $v_5 = 5^2 + 1 = 26$
    \end{multihideW}
  \end{minipage}
  \hfill
  \begin{minipage}[t]{0.31\linewidth}
    $w_n = \dfrac{120}{n + 1}$
    \begin{multihideW}[1em]
      $w_0 = 120 \div (0 + 1) = 120$

      $w_5 = 120 \div (5 + 1) = 20$
    \end{multihideW}
  \end{minipage}
\end{examples}

\subsection{Relation de récurrence}

\begin{example}
  Dans l'activité 1, le budget des repas vérifie $v_0 = 40\,000$ et $v_{n+1} = 1{,}1 \times v_n$.

  Pour calculer $v_2$, on calcule d'abord $v_1 = \hideM{1{,}1 \times 40\,000 = 44\,000}$, puis $v_2 = \hideM{1{,}1 \times 44\,000 = 48\,400}$.
\end{example}

\begin{definition}{Relation de récurrence}{}
  \begin{multihide}
    Une suite est définie par une \textbf{relation de récurrence} quand on donne son \textbf{premier terme} (par exemple $u_0$) et une relation qui donne $u_{n+1}$ \textbf{à partir du terme précédent} $u_n$.
  \end{multihide}
\end{definition}

\begin{method}{Calculer des termes avec une relation de récurrence}{}
  On calcule les termes \textbf{de proche en proche} : $u_1$ à partir de $u_0$, puis $u_2$ à partir de $u_1$, etc.
\end{method}

\begin{examples}
  Calculer les termes de rang $1$, $2$ et $3$ de chaque suite.

  \medskip
  \begin{minipage}[t]{0.48\linewidth}
    $u_0 = 3$ et $u_{n+1} = 2 \times u_n - 1$.
    \begin{multihideW}[1em]
      $u_1 = 2 \times u_0 - 1 = 2 \times 3 - 1 = 5$

      $u_2 = 2 \times u_1 - 1 = 2 \times 5 - 1 = 9$

      $u_3 = 2 \times u_2 - 1 = 2 \times 9 - 1 = 17$
    \end{multihideW}
  \end{minipage}
  \hfill
  \begin{minipage}[t]{0.48\linewidth}
    $a_0 = 10$ et $a_{n+1} = a_n + 4$.
    \begin{multihideW}[1em]
      $a_1 = a_0 + 4 = 10 + 4 = 14$

      $a_2 = a_1 + 4 = 14 + 4 = 18$

      $a_3 = a_2 + 4 = 18 + 4 = 22$
    \end{multihideW}
  \end{minipage}
\end{examples}

\subsection{Comparer les deux modes}

\renewcommand{\arraystretch}{1.3}
\begin{center}
\begin{tabularx}{\linewidth}{|L{3.6cm}|Y|Y|}
  \hline
  & \textbf{Formule explicite} & \textbf{Relation de récurrence} \\
  \hline
  Exemple & $u_n = 20 + 6n$ & $v_0 = 40\,000$ et $v_{n+1} = 1{,}1 \times v_n$ \\
  \hline
  On utilise\dots & \hide{son rang $n$} & \hide{le terme précédent} \\
  \hline
  Pour calculer $u_{10}$\dots & \hide{directement} & \hide{il faut $u_1$, $u_2$, \dots, $u_9$} \\
  \hline
\end{tabularx}
\end{center}

\begin{examples}
  Indiquer si chaque suite est définie par une formule explicite ou par une relation de récurrence.

  \renewcommand{\arraystretch}{1.3}
  \begin{center}
  \begin{tabularx}{\linewidth}{|c|Y|c|Y|}
    \hline
    $u_n = 5n - 1$ & \hide{explicite} & $v_0 = 2$ et $v_{n+1} = v_n - 3$ & \hide{récurrence} \\
    \hline
    $w_n = n^2$ & \hide{explicite} & $t_0 = 800$ et $t_{n+1} = 0{,}5 \times t_n$ & \hide{récurrence} \\
    \hline
  \end{tabularx}
  \end{center}
\end{examples}

\newpage
\section{Représentation graphique}

\begin{definition}{Représentation graphique d'une suite}{}
  \begin{multihide}
    Dans un repère, la représentation graphique d'une suite $u$ est le \textbf{nuage de points} de coordonnées $(n~;~u_n)$. On \textbf{ne relie pas} les points, car $n$ est un nombre entier.
  \end{multihide}
\end{definition}

\begin{method}{Représenter une suite}{}
  \begin{enumerate}
    \item On calcule les premiers termes et on les range dans un tableau.
    \item On place les points $(n~;~u_n)$ : le rang $n$ en abscisse, le terme $u_n$ en ordonnée.
  \end{enumerate}
\end{method}

\begin{example}
  \begin{minipage}[t]{0.5\linewidth}
    \vspace{0pt}
    Représenter la suite définie par $u_n = n^2 - 4n + 5$ pour $n$ allant de $0$ à $5$.

    \renewcommand{\arraystretch}{1.4}
    \begin{center}
    \begin{tabular}{|c|C{0.6cm}|C{0.6cm}|C{0.6cm}|C{0.6cm}|C{0.6cm}|C{0.6cm}|}
      \hline
      $n$ & $0$ & $1$ & $2$ & $3$ & $4$ & $5$ \\
      \hline
      $u_n$ & \hideM{5} & \hideM{2} & \hideM{1} & \hideM{2} & \hideM{5} & \hideM{10} \\
      \hline
    \end{tabular}
    \end{center}

    Calcul : $u_1 = \hideM{1^2 - 4 \times 1 + 5 = 2}$.

    \medskip
    \textbf{Remarque :} ici, les points ne sont pas alignés. Si $u_n = an + b$, les points sont \hide{alignés} sur la droite d'équation $y = ax + b$ (chapitre 3).
  \end{minipage}
  \hfill
  \begin{minipage}[t]{0.45\linewidth}
    \vspace{0pt}
    \centering
    \begin{tikzpicture}[x=0.8cm, y=0.34cm]
      \repere{0}{6}{0}{11}
      \hidetk{
        \foreach \n/\u in {0/5, 1/2, 2/1, 3/2, 4/5, 5/10} \fill[red] (\n,\u) circle (3pt);
      }
    \end{tikzpicture}
  \end{minipage}
\end{example}

\begin{method}{Lire un terme sur un graphique}{}
  Pour lire $u_n$, on part de l'abscisse $n$, on monte jusqu'au point, puis on lit son ordonnée.
\end{method}

\begin{example}
  \begin{minipage}[t]{0.45\linewidth}
    \vspace{0pt}
    \centering
    \begin{tikzpicture}[x=0.8cm, y=0.25cm]
      \repere{0}{6}{0}{16}
      \foreach \n/\u in {0/15, 1/12, 2/10, 3/9, 4/8.5, 5/8} \fill[blue] (\n,\u) circle (3pt);
    \end{tikzpicture}
  \end{minipage}
  \hfill
  \begin{minipage}[t]{0.5\linewidth}
    \vspace{1em}
    On a représenté une suite $v$ ci-contre.
    \begin{itemize}
      \item Le premier terme est $v_0 = \hideM{15}$.
      \item $v_2 = \hideM{10}$ et $v_3 = \hideM{9}$.
      \item Le rang du terme égal à $12$ est \hideM{1}.
    \end{itemize}
  \end{minipage}
\end{example}

\newpage
\section{Modéliser une situation par une suite}

\begin{method}{Traduire une phrase par une relation de récurrence}{}
  Chaque année (mois, jour\dots), la quantité\dots

  \smallskip
  \begin{tabularx}{\linewidth}{|Y|C{3.3cm}|Y|C{4.6cm}|}
    \hline
    \rule[-1.2em]{0pt}{3.2em}augmente de $b$ & $u_{n+1} = \hideM{u_n + b}$ & augmente de $t\,\%$ & $u_{n+1} = \hideM{\left(1 + \frac{t}{100}\right) \times u_n}$ \\
    \hline
    \rule[-1.2em]{0pt}{3.2em}diminue de $b$ & $u_{n+1} = \hideM{u_n - b}$ & diminue de $t\,\%$ & $u_{n+1} = \hideM{\left(1 - \frac{t}{100}\right) \times u_n}$ \\
    \hline
  \end{tabularx}
\end{method}

\begin{examples}
  \begin{enumerate}
    \setlength{\itemsep}{0.7em}
    \item Un EHPAD compte $80$ résidents en 2026. Chaque année, il accueille $5$ résidents de plus.
    On note $r_n$ le nombre de résidents en $2026 + n$.

    \smallskip
    \begin{multihideW}[1em]
      $r_0 = 80$ et $r_{n+1} = r_n + 5$.
    \end{multihideW}

    \item Une colonie de $500$ bactéries augmente de $20\,\%$ chaque heure.
    On note $b_n$ le nombre de bactéries au bout de $n$ heures. Calculer $b_1$ et $b_2$.

    \smallskip
    \begin{multihideW}[1em]
      Augmenter de $20\,\%$, c'est multiplier par $1{,}2$. Donc $b_0 = 500$ et $b_{n+1} = 1{,}2 \times b_n$.

      $b_1 = 1{,}2 \times 500 = 600$ et $b_2 = 1{,}2 \times 600 = 720$.
    \end{multihideW}

    \item Un patient reçoit $10$~mL d'un médicament. Chaque jour, son organisme en élimine $30\,\%$, puis on lui injecte $2$~mL.
    On note $q_n$ la quantité de médicament, en mL, au bout de $n$ jours. Calculer $q_1$, $q_2$ et $q_3$.

    \smallskip
    \begin{multihideW}[1em]
      Éliminer $30\,\%$, c'est multiplier par $0{,}7$ ; puis on ajoute $2$. Donc $q_0 = 10$ et $q_{n+1} = 0{,}7 \times q_n + 2$.

      $q_1 = 0{,}7 \times 10 + 2 = 9$ \qquad $q_2 = 0{,}7 \times 9 + 2 = 8{,}3$ \qquad $q_3 = 0{,}7 \times 8{,}3 + 2 = 7{,}81$
    \end{multihideW}
  \end{enumerate}
\end{examples}

\newpage
\section{Sens de variation}

\begin{definition}{}{}
  \begin{multihide}
    \begin{itemize}
      \item Une suite $u$ est \textbf{croissante} si $u_{n+1} \geqslant u_n$ pour tout $n \in \mathbb{N}$ : les termes augmentent.
      \item Une suite $u$ est \textbf{décroissante} si $u_{n+1} \leqslant u_n$ pour tout $n \in \mathbb{N}$ : les termes diminuent.
    \end{itemize}
    Certaines suites ne sont ni croissantes, ni décroissantes.
  \end{multihide}
\end{definition}

\begin{center}
  \begin{tabular}{C{0.31\linewidth}C{0.31\linewidth}C{0.31\linewidth}}
    \hide{croissante} & \hide{décroissante} & \hide{ni l'une ni l'autre} \\[0.3em]
    \begin{tikzpicture}[scale=0.45]
      \draw[->] (-0.3,0) -- (5.5,0); \draw[->] (0,-0.3) -- (0,5);
      \foreach \n/\u in {0/0.5, 1/1, 2/1.8, 3/2.8, 4/4.2} \fill[blue] (\n,\u) circle (4pt);
    \end{tikzpicture}
    &
    \begin{tikzpicture}[scale=0.45]
      \draw[->] (-0.3,0) -- (5.5,0); \draw[->] (0,-0.3) -- (0,5);
      \foreach \n/\u in {0/4.5, 1/3.2, 2/2.4, 3/1.9, 4/1.6} \fill[blue] (\n,\u) circle (4pt);
    \end{tikzpicture}
    &
    \begin{tikzpicture}[scale=0.45]
      \draw[->] (-0.3,0) -- (5.5,0); \draw[->] (0,-0.3) -- (0,5);
      \foreach \n/\u in {0/2, 1/3.5, 2/1.5, 3/4, 4/2.5} \fill[blue] (\n,\u) circle (4pt);
    \end{tikzpicture}
  \end{tabular}
\end{center}

\begin{method}{Conjecturer le sens de variation}{}
  On regarde un tableau de valeurs ou le nuage de points : si les termes augmentent, on conjecture que la suite est croissante ; s'ils diminuent, qu'elle est décroissante.
\end{method}

\begin{example}
  La suite $u_n = n^2 - 4n + 5$ de la partie II prend les valeurs $5$, $2$, $1$, $2$, $5$, $10$.

  Elle n'est \hide{ni croissante, ni décroissante}.

  Elle \hide{diminue} jusqu'au rang $2$, puis elle \hide{augmente}.
\end{example}

\begin{proposition}{}{}
  \begin{multihide}
    \begin{itemize}
      \item Si $u_{n+1} - u_n \geqslant 0$ pour tout entier naturel $n$, alors la suite $u$ est \textbf{croissante}.
      \item Si $u_{n+1} - u_n \leqslant 0$ pour tout entier naturel $n$, alors la suite $u$ est \textbf{décroissante}.
    \end{itemize}
  \end{multihide}
\end{proposition}

\begin{method}{Étudier le sens de variation par le calcul}{}
  On calcule la différence $u_{n+1} - u_n$, on la simplifie, on étudie son signe, puis on conclut.
\end{method}

\begin{examples}
  Étudier le sens de variation de chaque suite.

  \medskip
  \begin{minipage}[t]{0.42\linewidth}
    \textbf{1.} $u_0 = 50$ et $u_{n+1} = u_n - 3$.
    \begin{multihideW}[1em]
      $u_{n+1} - u_n = (u_n - 3) - u_n = -3$.

      $-3 < 0$, donc la suite $u$ est décroissante.
    \end{multihideW}
  \end{minipage}
  \hfill
  \begin{minipage}[t]{0.55\linewidth}
    \textbf{2.} $v_n = 2n + 7$.
    \begin{multihideW}[1em]
      $v_{n+1} = 2(n + 1) + 7 = 2n + 2 + 7 = 2n + 9$.

      $v_{n+1} - v_n = (2n + 9) - (2n + 7)$

      \hspace{4.4em}$= 2n + 9 - 2n - 7 = 2$.

      $2 > 0$, donc la suite $v$ est croissante.
    \end{multihideW}
  \end{minipage}
\end{examples}

\newpage
\section{Suites, tableur et Python}

\subsection{Tableur}

\begin{example}
  \begin{minipage}[t]{0.4\linewidth}
    \vspace{0pt}
    \renewcommand{\arraystretch}{1.2}
    \begin{tabular}{|>{\columncolor{gray!20}}c|C{1cm}|C{1.6cm}|C{1.8cm}|}
      \hline
      \rowcolor{gray!20} & A & B & C \\
      \hline
      1 & $n$ & $u_n$ & $v_n$ \\
      \hline
      2 & 0 & 20 & 40\,000 \\
      \hline
      3 & 1 & 26 & 44\,000 \\
      \hline
      4 & 2 & 32 & 48\,400 \\
      \hline
      5 & 3 & 38 & 53\,240 \\
      \hline
      6 & 4 & 44 & 58\,564 \\
      \hline
      7 & 5 & 50 & 64\,420,4 \\
      \hline
    \end{tabular}
  \end{minipage}
  \hfill
  \begin{minipage}[t]{0.57\linewidth}
    \vspace{0pt}
    On reprend les suites de l'activité 1 : $u_n = 20 + 6n$, et $v_0 = 40\,000$, $v_{n+1} = 1{,}1 \times v_n$.

    \medskip
    \begin{itemize}
      \item Formule explicite : en B2, on saisit \hide{\texttt{=20+6*A2}}, puis on recopie vers le bas.
      \item Récurrence : en C2, on saisit $40\,000$. En C3, on saisit \hide{\texttt{=1,1*C2}}, puis on recopie vers le bas.
      \item En recopiant, la formule de C3 devient en C4 : \hide{\texttt{=1,1*C3}}.
    \end{itemize}
  \end{minipage}
\end{example}

\begin{method}{Déterminer le rang à partir duquel un seuil est dépassé}{}
  On calcule les termes un par un et on repère le \textbf{premier rang} $n$ où $u_n$ dépasse le seuil.
\end{method}

\begin{example}
  À partir de quelle année le budget des repas dépasse-t-il $60\,000$~€~?
  \begin{multihideW}[1em]
    D'après le tableur, $v_4 = 58\,564 < 60\,000$ et $v_5 = 64\,420{,}4 > 60\,000$.

    Le seuil est dépassé à partir du rang $n = 5$, c'est-à-dire en $2026 + 5 = 2031$.
  \end{multihideW}
\end{example}

\subsection{Python}

\begin{example}
  \begin{minipage}[c]{0.42\linewidth}
    \begin{python}
      def u(n):\\
      \tab return 20 + 6*n
    \end{python}
  \end{minipage}
  \hfill
  \begin{minipage}[c]{0.55\linewidth}
    \raggedright
    \textbf{Formule explicite.} Cette fonction renvoie $u_n = 20 + 6n$.

    L'instruction \texttt{u(10)} renvoie \hide{$80$}.
  \end{minipage}
\end{example}

\begin{example}
  \begin{minipage}[c]{0.42\linewidth}
    \begin{python}
      def v(n):\\
      \tab terme = 40000\\
      \tab for i in range(n):\\
      \tab\tab terme = 1.1*terme\\
      \tab return terme
    \end{python}
  \end{minipage}
  \hfill
  \begin{minipage}[c]{0.55\linewidth}
    \raggedright
    \textbf{Relation de récurrence.} Cette fonction renvoie $v_n$, avec $v_0 = 40\,000$ et $v_{n+1} = 1{,}1 \times v_n$.
    \begin{itemize}
      \item Au départ, \texttt{terme} contient \hide{$v_0 = 40\,000$}.
      \item La boucle répète $n$ fois : \hide{multiplier par $1{,}1$}.
      \item \texttt{v(2)} renvoie \hide{$48\,400$}.
    \end{itemize}
  \end{minipage}
\end{example}

\begin{example}
  \begin{minipage}[c]{0.42\linewidth}
    \begin{python}
      def q(n):\\
      \tab terme = \hide{10}\\
      \tab for i in range(n):\\
      \tab\tab terme = \hide{0.7*terme + 2}\\
      \tab return terme
    \end{python}
  \end{minipage}
  \hfill
  \begin{minipage}[c]{0.55\linewidth}
    \raggedright
    Compléter la fonction qui renvoie la quantité de médicament $q_n$, avec $q_0 = 10$ et $q_{n+1} = 0{,}7 \times q_n + 2$.

    \smallskip
    \textbf{Attention :} en Python, la virgule décimale s'écrit avec un \textbf{point} : $0{,}7$ s'écrit \texttt{0.7}.
  \end{minipage}
\end{example}

\vfill

\begin{savoir}{}{}
\begin{itemize}[label=--]
  \item Utiliser les notations $u(n)$ et $u_n$,
  \item Calculer un terme avec une formule explicite ou une relation de récurrence,
  \item Reconnaître le mode de génération d'une suite,
  \item Représenter une suite par un nuage de points ; lire un terme sur un graphique,
  \item Traduire une situation par une relation de récurrence,
  \item Conjecturer le sens de variation ; le démontrer avec le signe de $u_{n+1} - u_n$,
  \item Compléter une formule de tableur ; trouver le rang où un seuil est dépassé,
  \item Compléter un programme Python qui calcule un terme d'une suite.
\end{itemize}
\end{savoir}

\end{document}
